\documentclass[10pt,a4paper]{amsart}
\usepackage{amsmath,amssymb,amsthm,graphicx}
% Retained preamble for the Stacks Project chapter snapshot in this directory.
% The chapter's own preamble.tex is a copy of the shared Stacks Project preamble; the block
% below is sliced MECHANICALLY, line for line, from it. Nothing was added, removed or
% redefined; the wrapper class, the xr-hyper/hyperref wiring and the externaldocument list
% are omitted because the assembler supplies the document class.
%
% lines taken from preamble.tex: 3, 5, 9, 11, 12, 156, 157, 158, ... (51 lines total)
%
% package coverage restored from the chapter's own preamble (2026-10-13): xy, xy-2cell

\usepackage[all]{xy}
\xyoption{2cell}
\UseAllTwocells

\newenvironment{reference}{\comment}{\endcomment}
\newenvironment{slogan}{\comment}{\endcomment}
\newenvironment{history}{\comment}{\endcomment}
\theoremstyle{plain}
\newtheorem{theorem}[subsection]{Theorem}
\newtheorem{proposition}[subsection]{Proposition}
\newtheorem{lemma}[subsection]{Lemma}
\theoremstyle{definition}
\newtheorem{definition}[subsection]{Definition}
\newtheorem{example}[subsection]{Example}
\newtheorem{exercise}[subsection]{Exercise}
\newtheorem{situation}[subsection]{Situation}
\theoremstyle{remark}
\newtheorem{remark}[subsection]{Remark}
\newtheorem{remarks}[subsection]{Remarks}
\numberwithin{equation}{subsection}
\def\lim{\mathop{\mathrm{lim}}\nolimits}
\def\colim{\mathop{\mathrm{colim}}\nolimits}
\def\Spec{\mathop{\mathrm{Spec}}}
\def\Hom{\mathop{\mathrm{Hom}}\nolimits}
\def\Ext{\mathop{\mathrm{Ext}}\nolimits}
\def\SheafHom{\mathop{\mathcal{H}\!\mathit{om}}\nolimits}
\def\SheafExt{\mathop{\mathcal{E}\!\mathit{xt}}\nolimits}
\def\Sch{\mathit{Sch}}
\def\Mor{\mathop{\mathrm{Mor}}\nolimits}
\def\Ob{\mathop{\mathrm{Ob}}\nolimits}
\def\Sh{\mathop{\mathit{Sh}}\nolimits}
\def\NL{\mathop{N\!L}\nolimits}
\def\CH{\mathop{\mathrm{CH}}\nolimits}
\def\proetale{{pro\text{-}\acute{e}tale}}
\def\etale{{\acute{e}tale}}
\def\QCoh{\mathit{QCoh}}
\def\Ker{\mathop{\mathrm{Ker}}}
\def\Im{\mathop{\mathrm{Im}}}
\def\Coker{\mathop{\mathrm{Coker}}}
\def\Coim{\mathop{\mathrm{Coim}}}
\DeclareMathSymbol{\boxtimes}{\mathbin}{AMSa}{"02}
\def\QCohstack{\mathcal{QC}\!\mathit{oh}}
\def\Cohstack{\mathcal{C}\!\mathit{oh}}
\def\Spacesstack{\mathcal{S}\!\mathit{paces}}
\def\Quotfunctor{\mathrm{Quot}}
\def\Hilbfunctor{\mathrm{Hilb}}
\def\Curvesstack{\mathcal{C}\!\mathit{urves}}
\def\Polarizedstack{\mathcal{P}\!\mathit{olarized}}
\def\Complexesstack{\mathcal{C}\!\mathit{omplexes}}
\def\Pic{\mathop{\mathrm{Pic}}\nolimits}
\def\Picardstack{\mathcal{P}\!\mathit{ic}}
\def\Picardfunctor{\mathrm{Pic}}
\def\Deformationcategory{\mathcal{D}\!\mathit{ef}}
\usepackage{cleveref}
\crefname{originallemma}{original lemma}{original lemmas}
\Crefname{originallemma}{Original lemma}{Original lemmas}
\makeatletter
\newlabel{example-finite-projective-modules}{{0D3I}{source}{}{sourcepointer-1}{}}
\newlabel{example-finite-projective-modules@cref}{{[originallemma][0][]0D3I}{[1][0][]source}{}{}{}}
\makeatother
\crefname{originallemma}{original lemma}{original lemmas}
\Crefname{originallemma}{Original lemma}{Original lemmas}
\makeatletter
\newlabel{formal-defos-definition-infinitesimal-auts}{{06JQ}{source}{}{sourcepointer-2}{}}
\newlabel{formal-defos-definition-infinitesimal-auts@cref}{{[originallemma][0][]06JQ}{[1][0][]source}{}{}{}}
\makeatother
\title{Stacks, Examples of Deformations: lemma-finite-projective-modules-TI}
\author{The Stacks Project Authors}
\date{stacks-project:examples-defos.tex:local-snapshot-20261012; GNU FDL 1.2 or later}
\begin{document}
\maketitle
\noindent\textit{Editorial note.} Statement span is lines 429-439 of sources/stacks-examples-defos/pinned-original-source.tex; proof span is lines 441-473, the complete author proof (31 body lines).\par\medskip
\noindent\textit{Editorial note.} Navigation stubs: 2 cited-but-not-included labels are supplied as non-load-bearing pointers with their pinned tags where the snapshot tag index has them.\par\medskip
\noindent\textit{Editorial note.} extra\_packages is empty on purpose: the retained preamble already loads xy with options and a second declaration clashes in the assembler wrapper.\par\medskip
\noindent\textit{Editorial note.} Text is retained verbatim, including the author's own typos and placeholders.\par\medskip
\noindent\textit{Editorial note.} Queue index 15 of queue\_examples-defos.json (pb=1471, stmt 429-439, proof 441-473).\par\medskip
\noindent\textit{Allowed standard prerequisites (editorial).}\begin{itemize}
\item Explicit prerequisite boundary. The cited inputs are prerequisites; the argument in the retained span is not replaced by any cited result.
\end{itemize}\medskip
\ifdefined\hypertarget\hypertarget{sourcepointer-1}{}\fi
\noindent\textit{Source reference.} Original lemma 0D3I: sources/stacks-examples-defos/pinned-original-source.tex, source lines 359--369. Cited by the retained statement or proof but not part of the collected spans; supplied only as a navigation pointer. The argument does not depend on its content.\par\medskip
\ifdefined\hypertarget\hypertarget{sourcepointer-2}{}\fi
\noindent\textit{Source reference.} Original lemma 06JQ: sources/stacks-examples-defos/pinned-original-source.tex, source lines 473--473. Cited by the retained statement or proof but not part of the collected spans; supplied only as a navigation pointer. The argument does not depend on its content.\par\medskip
\begin{lemma}
\label{lemma-finite-projective-modules-TI}
In Example \ref{example-finite-projective-modules}
let $V$ be a finite dimensional $k$-vector space. Then
$$
T\Deformationcategory_V = (0)
\quad\text{and}\quad
\text{Inf}(\Deformationcategory_V) = \text{End}_k(V)
$$
are finite dimensional.
\end{lemma}

\begin{proof}
With $\mathcal{F}$ as in Example \ref{example-finite-projective-modules}
set $x_0 = (k, V) \in \Ob(\mathcal{F}(k))$.
Recall that $T\Deformationcategory_V = T_{x_0}\mathcal{F}$
is the set of isomorphism
classes of pairs $(x, \alpha)$ consisting of an object $x$ of $\mathcal{F}
$ over the dual numbers $k[\epsilon]$ and a morphism
$\alpha : x \to x_0$ of $\mathcal{F}$ lying over $k[\epsilon] \to k$.

\medskip\noindent
Up to isomorphism, there is a unique pair $(M, \alpha)$ consisting of a
finite projective module $M$ over $k[\epsilon]$
and $k[\epsilon]$-linear map $\alpha : M \to V$
which induces an isomorphism $M \otimes_{k[\epsilon]} k \to V$.
For example, if $V = k^{\oplus n}$, then we take
$M = k[\epsilon]^{\oplus n}$ with the obvious map $\alpha$.

\medskip\noindent
Similarly, $\text{Inf}(\Deformationcategory_V) = \text{Inf}_{x_0}(\mathcal{F})$
is the set of automorphisms
of the trivial deformation $x'_0$ of $x_0$ over $k[\epsilon]$.
See Formal Deformation Theory, Definition
\ref{formal-defos-definition-infinitesimal-auts} for details.

\medskip\noindent
Given $(M, \alpha)$ as in the second paragraph, we see that an element of
$\text{Inf}_{x_0}(\mathcal{F})$ is an automorphism $\gamma : M \to M$ with
$\gamma \bmod \epsilon = \text{id}$. Then we can write
$\gamma = \text{id}_M + \epsilon \psi$ where
$\psi : M/\epsilon M \to M/\epsilon M$ is $k$-linear.
Using $\alpha$ we can think of $\psi$ as an element of
$\text{End}_k(V)$ and this finishes the proof.
\end{proof}

\end{document}
